\documentclass[11pt]{article}
    \title{\textbf{Computer Architecture}\\[0.5em]
    \large Project Step 3: ALU Control Unit}
    
    \author{Group ``Orange": Chloe MacGregor, Nathaniel Jefferson, Rex Almerico}
    \date{April 29, 2026}

    \addtolength{\topmargin}{-3cm}
    \addtolength{\textheight}{3cm}

\usepackage{listings}
\usepackage{color}

\usepackage{graphicx}
\graphicspath{ {images/} }

\definecolor{dkgreen}{rgb}{0,0.6,0}
\definecolor{gray}{rgb}{0.5,0.5,0.5}
\definecolor{mauve}{rgb}{0.58,0,0.82}

\lstset{frame=tb,
  language=Verilog,
  aboveskip=3mm,
  belowskip=3mm,
  showstringspaces=false,
  columns=flexible,
  basicstyle={\small\ttfamily},
  numbers=none,
  numberstyle=\tiny\color{gray},
  keywordstyle=\color{blue},
  commentstyle=\color{dkgreen},
  stringstyle=\color{mauve},
  breaklines=true,
  breakatwhitespace=true,
  tabsize=3
}

\begin{document}

\maketitle
\thispagestyle{empty}

\section{Objectives}
The goal of this step in the project is to integrate the 4-bit ALU components from Step 2 into a control unit that selects the appropriate operation based on a 4-bit opcode. This creates a complete ALU that can perform arithmetic and logic operations on 4-bit inputs.

\section{ALU Control Unit Design}
The control unit takes a 4-bit opcode and two 4-bit inputs A and B, and produces an 8-bit output based on the selected operation. It instantiates all the 4-bit modules from Step 2 and uses a case statement to select the output.

\subsection{Operation Codes}
The following table lists the supported operations and their corresponding opcodes:

\begin{itemize}
\item 4'b0000 (0)  - ADD   : out[4:0]  = \{Cout, A+B\}
\item 4'b0001 (1)  - SUB   : out[4:0]  = \{Bout, A-B\}
\item 4'b0010 (2)  - MUL   : out[7:0]  = \{upper4, lower4\} of A*B
\item 4'b0011 (3)  - DIV   : out[7:0]  = \{remainder, quotient\} of A/B
\item 4'b0100 (4)  - AND   : out[3:0]  = A \& B
\item 4'b0101 (5)  - OR    : out[3:0]  = A | B
\item 4'b0110 (6)  - XOR   : out[3:0]  = A \^{} B
\item 4'b0111 (7)  - SHIFT : out[7:4]  = bits shifted out, out[3:0] = shifted result \\
                          B[3]=dir(0=L,1=R), B[2:1]=amount(0-3), B[0]=fill
\item 4'b1100 (12) - NAND  : out[3:0]  = \~(A \& B)
\item 4'b1101 (13) - NOR   : out[3:0]  = \~(A | B)
\item 4'b1110 (14) - XNOR  : out[3:0]  = \~(A \^{} B)
\item 4'b1111 (15) - NOT   : out[7:4]  = \~B,  out[3:0] = \~A
\end{itemize}

\subsection{Verilog Code}
\lstinputlisting[breaklines]{step-3/control.v}

\pagebreak

\section{Unit Tests}
The control unit is tested with a comprehensive testbench that verifies all operations with various inputs.

\lstinputlisting[breaklines]{step-3/tests/test_control.v}

\pagebreak

\section{Waveforms for Each Operation}

\subsection{ADD (Opcode 0000)}
\includegraphics[width=\textwidth]{images/control_ADD_wvf.png}

\subsection{SUB (Opcode 0001)}
\includegraphics[width=\textwidth]{images/control_SUB_wvf.png}

\subsection{MUL (Opcode 0010)}
\includegraphics[width=\textwidth]{images/control_MUL_wvf.png}

\subsection{DIV (Opcode 0011)}
\includegraphics[width=\textwidth]{images/control_DIV_wvf.png}

\subsection{AND (Opcode 0100)}
\includegraphics[width=\textwidth]{images/control_AND_wvf.png}

\subsection{OR (Opcode 0101)}
\includegraphics[width=\textwidth]{images/control_OR_wvf.png}

\subsection{XOR (Opcode 0110)}
\includegraphics[width=\textwidth]{images/control_XOR_wvf.png}

\subsection{SHIFT (Opcode 0111)}
\includegraphics[width=\textwidth]{images/control_SHIFT_wvf.png}

\subsection{NAND (Opcode 1100)}
\includegraphics[width=\textwidth]{images/control_NAND_wvf.png}

\subsection{NOR (Opcode 1101)}
\includegraphics[width=\textwidth]{images/control_NOR_wvf.png}

\subsection{XNOR (Opcode 1110)}
\includegraphics[width=\textwidth]{images/control_XNOR_wvf.png}

\subsection{NOT (Opcode 1111)}
\includegraphics[width=\textwidth]{images/control_NOT_wvf.png}

\section{Conclusion}
In this step of the project, we successfully integrated the individual 4-bit ALU components into a unified control unit. The control unit uses opcodes to select between different arithmetic and logic operations, producing the correct 8-bit output. This completes the ALU design for this computer architecture project.

\end{document}